\scanpage{5}
\textcircled{$6^\circ$}. Utilisant la condition \textcircled{$7^\circ$}, on est ramené à montrer que $X_i\xrightarrow{u_i}F$ est modulaire en le centre $x_i$ de $X_i$. Comme la restriction $u_0:X_0\mathrel{\unsure{-|\longrightarrow}}F_0$ est modulaire (donc \deleted{\unsure{le}} $k(x_i)$ est un corps de définition pour $u_i:k(x_i)\longrightarrow F$ sur $S$] il reste à vérifier la condition modulaire. Nous allons \deleted{\unsure{commencer}} d'abord démontrer le théorème dans le cas où la condition $7^\circ$ est remplie, en y exigeant seulement $\unsure{\eta|y}$ \uline{pro-modulaire} \inserted{sur $S$} en $z$. De plus, $X$ est loc. qu. fini sur $S$ par construction, et séparé sur $S$ par $6^\circ$.
\margin{Dégager un lemme à cet effet}
\margin{Pour l'instant on n'utilise \deleted{\unsure{pas}} \inserted{\unsure{que la}} \deleted{\unsure{la plus}} \inserted{\unsure{forme}} \unsure{affaiblie} \deleted{\unsure{nouvelle}} de $7^\circ$}

\textcircled{3} $\dim S=n>0$, $S$ local. On sait par $2)$ que $F_{\hat S}$ est représentable par un $X_{\unsure{\hat S}}$ \inserted{loc.} quasi-fini, et séparé sur $S$. Il en résulte que la donnée de descente \unsure{canon.} sur $F_{\hat S}$ est effective, \deleted{Utilisant} d'où un $X$ sur $S$, \inserted{loc.} quasi-fini sur $S$ et séparé sur $S$.

Utilisant $1^\circ$, b) on trouve qu'il représente $F$.

\textcircled{4} $\dim S=n$, $S$ quelconque. \deleted{\unsure{Comme le}} On utilise le

\uline{Lemme} \inserted{$F$ comme dessus.} Supposons que pour $\forall\ a\in S$, $F_{S^a}$ sur $S^a=\operatorname{Spec}\mathcal O_{S,a}$ soit représentable. Alors $F$ est représentable.
\margin{\figure{\begin{tikzpicture}\draw (0,0)--(0,1.4);\end{tikzpicture}}}

On utilise la prop $3$ des notes jointes.
\mspage{4}
